\section{A Design Space}
\label{sec:mnemos}

\Cref{thm:main} gives conditions a storage layout can satisfy, and \cref{prop:capacity,prop:derived}
give two it cannot satisfy by layout alone. This section describes three designs that sit at
different points in that space. All three are implemented in Rust with a Python binding used only
by the evaluation harness and the embedding model; \cref{sec:eval} measures them, and
\cref{sec:discussion} states what each does not achieve.

\begin{itemize}
\item \textbf{Mnemos-RS:} a record store plus a separate index: the natural padded design, and the one a
reader would write first.
\item \textbf{Mnemos-P:} one sealed partition file of $B$ equal regions, placed by content. Intended for
stores larger than memory.
\item \textbf{Mnemos-M:} the same sealed file, placed at random, searched entirely in memory, and flushed
on a clock. Intended for stores that fit in memory, which for agent memory is the common case.
\end{itemize}

\Cref{tab:crossscheme} in \cref{sec:discussion} summarizes which channel each closes, beside the
alternatives; the rest of this section says how.

\begin{figure*}[t]
\centering
\includegraphics[width=\linewidth]{figures/fig_architecture.pdf}
\caption{Where each design puts a record and where it searches. Both write whole sealed regions of
one fixed-size file, so neither reveals a region's occupancy. They differ in two places, and both
differences are the security argument: Mnemos-P sends a record to the region its \emph{content}
selects and reads $s\,k_{\mathrm{probe}}$ regions back from disk to answer a query, while Mnemos-M
draws the region uniformly and answers queries entirely from the locked arena, touching no disk at
all.}
\label{fig:architecture}
\end{figure*}

\subsection{The Shared Layout}
\label{sec:mnemos-layout}

Mnemos-P and Mnemos-M share a file of $B$ equal regions of $C$ slots. A slot holds two
authenticated ciphertexts: a fixed-size header carrying the record identifier and its quantized
embedding, and a fixed-size payload of one or more blocks carrying the text and metadata. Every
ciphertext uses XChaCha20-Poly1305 under a fresh 192-bit random nonce, whose length makes a
counter unnecessary and so keeps no nonce state on disk; a free slot is uniformly random bytes of
the slot length. A single authenticated cipher is used throughout. This satisfies S1 and S2: the
file lengths are functions of $B$, $C$, the embedding dimension and the block size, and no slot's
bytes distinguish a live record from free space without the key.

Both designs write by rewriting \emph{whole regions}, so a write reveals which region changed and
never which slot within it, and hence never a region's occupancy. Both flush exactly $w+1$ regions.
What differs is how a region is chosen, and that difference is the whole of the security
difference between them.

\subsection{Mnemos-P: Placement by Content}
\label{sec:mnemos-p}
\label{sec:mnemos-write}

A record's region is the keyed permutation of its nearest centroid's partition. The permutation
exists so that an adversary holding a public partition table, one trained on public text, cannot
map a topic to a region directly. A secret \emph{rotation} of the centroids would not achieve
this, because a rotation preserves inner products and therefore preserves the nearest-centroid
assignment; permuting physical positions is what hides the correspondence.

The permutation is weaker than it appears against the RW class. Genuine probes are the
$k_{\mathrm{probe}}$ nearest partitions to a query, so the regions read together are geometric
neighbours, and the co-occurrence graph of observed reads can be matched against the graph the
public table predicts. \Cref{sec:eval-permutation} recovers most of the permutation from search
traffic alone when probes are not randomized, and part of it at the shipped spread. Against the W
class, which sees no reads, the permutation holds. This is stated here rather than in the
limitations because it changes what the mechanism is for: it protects against a write observer
with a public table, not against a read observer.

A record whose partition is full is refused rather than spilled, so that no history can make a
write land in a partition it would not otherwise reach. By \cref{prop:capacity} a refusal that
wrote nothing would itself be the leak, so a refused write performs the same region rewrite an
accepted one would have: the store owes the intended region its rewrite whether or not the record
was admitted. \Cref{sec:eval-refusal} measures both arms.

Two further measures address the multi-snapshot channel of \cref{sec:artifacts}. The whole file is
re-encrypted when the store is locked, so any version diff spanning a session boundary covers every
region and carries nothing; and the maintenance cursor's starting position is drawn afresh at each
unlock, so its resting position no longer counts the flushes performed before $\tau$.

\subsection{Mnemos-M: Placement at Random, Search in Memory}
\label{sec:mnemos-m}

Assumption A3 already grants that the adversary cannot read the unlocked process's memory, and in
the local setting there is no remote server that must hold the index. Mnemos-M takes both facts at
face value.

\paragraph{Placement} A record's slot is drawn uniformly from the store's free slots, at the moment
it is accepted, before its content has influenced anything on disk. Its region is that slot's
region. The $w$ record-carrying rewrites of a flush and the padding rewrites are therefore draws
from the \emph{same} distribution, which is what Mnemos-P cannot say: there, a real record's region
is a function of its content and only the padding is uniform. A version diff, at any granularity,
reveals which regions changed and nothing about what they hold.

\paragraph{Search} At unlock the store is read once, sequentially, and every header, and by default
every payload, is decrypted into the locked arena. A search ranks the resident headers there and
returns results from resident payloads, so it performs \emph{no disk access at all}. The search,
result and query-equality channels are not mitigated but absent. The cost is memory, which
\cref{sec:eval-perf} reports against $n$.

\paragraph{Logical partitions} Exhaustive ranking is linear in $n$, and \cref{sec:eval-perf}
finds the partitioned design overtaking it on search between \REtenCrossLo{} and \REtenCrossHi{}
records.
The remedy is to notice what the partition table actually costs. It cost Mnemos-P its write
channel because Mnemos-P's partitions are \emph{physical}: a record's region on disk is its
partition. Nothing requires that. Mnemos-M keeps each resident record's nearest-centroid
partition in the locked arena beside its header and, when configured to, ranks only the records
of the query's $k$ nearest partitions. The disk layout is unchanged, uniformly random, and the
probing happens in memory, so the index changes nothing either adversary class observes; the
auditor confirms that such a search still writes no byte and emits no event. What it trades is
recall, which \cref{sec:eval-recall} measures, and time at unlock, when every resident record's
partition is recomputed, which \cref{sec:eval-perf} measures. Separating the logical index from the
physical layout is what lets one store search sublinearly without placing records by content.

\paragraph{Flushes} A flush occurs every $\Delta$ while the store is unlocked, whether or not
anything was written. This is what closes the derived-write channel of \cref{prop:derived}: the
occurrence of a flush no longer implies that the agent wrote, so an adversary that supplies a
document and watches for a write learns nothing, and the deduplication attack of
\cref{sec:eval-dedup} fails against it. It also bounds deletion in wall-clock time rather than in
flushes, and hides write timing up to $\Delta$.

\paragraph{Capacity} There is no per-partition capacity, so \cref{prop:capacity} does not apply.
The only bound is global, and global occupancy is a function of the record count, which the game
already equalizes. The practical consequence is large: \cref{sec:eval-refusal} shows Mnemos-P
refusing its first write at roughly a quarter of nominal capacity under a skewed corpus, while
Mnemos-M admits every slot.

\subsection{Search and Probe Mechanisms}
\label{sec:mnemos-read}

Mnemos-M's search reads nothing, so this section concerns Mnemos-P. A search reads whole regions
and nothing else: it reads the regions of its probe set, decrypts the headers they contain, ranks
the candidates in memory, and decrypts the $\mathit{top\_k}$ payloads from the region images it has
already read. No further, result-dependent access reaches the disk, which removes the channel by
which a returned cohort-H record could be linked to a query topic (\cref{sec:eval-search}).

Every query reads the same number of regions, chosen by one of two mechanisms: the default reads
the $k_{\mathrm{probe}}$ genuine nearest partitions plus $(s-1)k_{\mathrm{probe}}$ fresh random
ones; the alternative reads one of a set of fixed disjoint groups, so repeated queries are
recognised only at group granularity, with a deterministic anonymity set. Neither hides query
equality from an overlap-threshold adversary (\cref{sec:eval-query}); only the grouped mechanism
resists permutation recovery (\cref{sec:eval-permutation}).

\subsection{Deletion and Update}
\label{sec:mnemos-delete}

In both designs a deletion records a tombstone in the sealed manifest and writes nothing else at
the time of the call, so the operation does not name the deleted record's region. The
counter-selected region of each flush is scanned for tombstoned slots, which are freed as that
region is rewritten, so a tombstoned record is physically removed within $B$ flushes at the latest.
An update is a deletion followed by a fresh write and therefore also writes nothing at the time of
the call. This is what lets the post-$\tau$ program delete and update cohort-H records without
leaking their partitions.

For Mnemos-P the bound is in flushes, which for an agent that writes rarely can mean days;
Mnemos-M's clock makes it $B\Delta$. Neither reaches copies a synchronization provider already
retains.

\subsection{Key Custody}
\label{sec:mnemos-keys}

The master key is wrapped by a hardware key store, the Windows TPM through the CNG Platform Crypto
Provider in the current implementation, and is unwrapped into a locked, promptly zeroized memory
region only while the store is unlocked. All working keys are derived from it by HKDF. Decrypted
records and query vectors live only in the locked region, which for Mnemos-M is also where the
resident index lives.

Since page locking does not prevent hibernation, custody is paired with locking the store on
suspend and hibernate notifications (\cref{sec:artifacts,sec:eval-artifacts}).

\subsection{Instrumentation}
\label{sec:mnemos-events}

Every disk-visible operation emits one typed event naming the file regions it touched, and that
stream is the transcript the attack harness consumes, so an attack cannot structurally observe more
than the storage adversary does. \Cref{sec:eval-audit} describes the separate byte-level audit that
checks the stream against the disk, and what it covers for each design.
